As machine learning systems are increasingly deployed in socially critical domains---including criminal justice, lending, healthcare, and hiring---ensuring that these systems do not discriminate based on legally protected attributes such as race, sex, and age has become a paramount concern. Existing fairness analysis tools predominantly operate as post-training evaluation frameworks, requiring practitioners to complete the full model development lifecycle before assessing bias. This report presents \textbf{FairLint-DL}, a Visual Studio Code extension that implements a \emph{shift-left} approach to fairness testing by enabling pre-training, IDE-native bias detection directly on tabular datasets.

FairLint-DL trains a configurable deep neural network as a proxy model and applies information-theoretic \emph{Quantitative Individual Discrimination} (QID) metrics---grounded in Shannon entropy and min-entropy---to quantify the causal influence of protected attributes on model predictions. The system implements a two-phase gradient-guided search algorithm for discovering discriminatory instances, a causal debugging pipeline that localizes bias to specific network layers and neurons via sensitivity analysis, and dual explainability engines using SHAP (SHapley Additive exPlanations) and LIME (Local Interpretable Model-agnostic Explanations) for feature-level attribution.

Evaluation on the Adult Census Income dataset (32,561 records, 14 features) reveals severe fairness concerns: 95.8\% of analyzed instances exhibit QID above the 0.1-bit significance threshold, with a mean QID of 0.640 bits. The disparate impact ratio of 0.581 violates the four-fifths legal standard, and group fairness metrics confirm systematic disparities across demographic parity, equalized odds, and equal opportunity for the protected attributes of sex, race, age, and national origin. Causal debugging identifies Layer~2 (32 neurons) as the most bias-sensitive layer, with Neuron~30 exhibiting the highest impact score of 1.362. The composite fairness score of approximately 41 out of 100 rates the dataset as ``Concerning.'' FairLint-DL produces these results within 12~seconds on cached models, demonstrating the feasibility of integrating comprehensive fairness analysis into the developer workflow without significant overhead.
